\section{Preliminary and Related Work}\label{sec:preliminary}

\subsection{Retrieval-augmented generation and HyDE-style query generation}
Retrieval-augmented generation conditions a generator on passages retrieved
from a non-parametric index, originally to inject up-to-date or
domain-specific knowledge into a frozen language model
\citep{lewis2020rag,izacard2021fid}. A practical limitation of the standard
recipe is that the literal user query is often a poor embedding target: it is
under-specified, missing domain terminology, or phrased differently from the
controlling source text. Query-side generation methods address this by
turning the query into a richer pseudo-document before retrieval. HyDE
synthesizes a hypothetical answer passage and embeds that passage instead of
the query \citep{gao2023hyde}; Query2Doc expands the query with a generated
document fragment \citep{wang2023query2doc}; GenRead replaces the retriever
with a generator that emits context documents \citep{yu2023genread}; and HyQE
inverts the direction by generating hypothetical queries for indexed passages
\citep{zhou2024hyqe}. \hyre{} sits in this family but is conditioned on a
\emph{snap answer}: the model first commits to a tentative legal answer to the
question, then writes the controlling rule passage in service of that
commitment. The embedding therefore reflects a proposed legal solution rather
than a generic pseudo-document, which biases retrieval toward authorities
that bear on the specific outcome at issue rather than toward topical
neighbors of the surface question.

\subsection{Adaptive and agentic RAG}
A growing literature treats retrieval as a controllable action rather than a
fixed pipeline stage. FLARE retrieves only when the generator's
next-token confidence dips \citep{jiang2023flare}; Self-RAG trains the
generator to emit reflection tokens that gate retrieve, critique, and rewrite
steps \citep{asai2024selfrag}; CRAG grades retrieved evidence and applies
corrective web search when the retrieved set is judged insufficient
\citep{yan2024crag}; Adaptive-RAG routes a question to a no-retrieval,
single-step, or multi-step pipeline based on a learned question-complexity
classifier \citep{jeong2024adaptive}; Adaptive HyDE thresholds the HyDE
retriever on per-query confidence to skip generation when retrieval already
suffices \citep{lei2025adaptivehyde}; Speculative RAG runs a small drafter
against a large verifier to amortize generation cost
\citep{wang2025specrag}; and Iter-RetGen \citep{shao2023iterretgen} together
with IRCoT \citep{trivedi2023ircot} interleave retrieval with intermediate
reasoning steps so that later queries condition on earlier ones. Our
\controller{} differs in two ways. First, we do not route by question
complexity or by trained reflection tokens; we route by \emph{bottleneck type
measured from a calibration trace} on each benchmark slice. Second, the
intervention space is \emph{heterogeneous} (rewrite, \hyre{}, state-filtered
retrieval, conservative verifier, disagreement arbitration, and
reject/escalate) rather than retrieve-depth or retrieve-or-not variants of
the same skeleton. Source-gated diagnostic measurement, with the bottleneck
labeling and intervention selection it implies, is the contribution.

\subsection{Legal NLP benchmarks and legal RAG}
Legal evaluation has matured around a small set of multi-task and reasoning
benchmarks: LegalBench \citep{guha2023legalbench}, which includes the SCALR
holdings task we evaluate; CaseHOLD \citep{zheng2021casehold}, the
five-option holding-completion benchmark; and the reasoning-focused legal
retrieval benchmark introducing BarExam and HousingQA
\citep{zheng2025legalretrieval}. These are paired with large legal corpora
such as LexGLUE \citep{chalkidis2022lexglue} and Pile of Law
\citep{henderson2022pileoflaw}. Recent work targets retrieval-aware
evaluation directly: LegalBench-RAG factors out retrieval from generation on
LegalBench prompts \citep{pipitone2024legalbenchrag}, and Legal RAG Bench
extends the protocol across additional legal tasks
\citep{butler2026legalragbench}. The reliability picture is more sobering:
the Stanford hallucination audit of commercial legal AI products documents
substantial fabrication and miscitation \citep{magesh2024hallucinations},
showing that legal RAG cannot be summarized by a single accuracy number; and
\citet{vaddi2026smallmodels} reports that adding legal retrieval helps or
hurts depending on the task. Our work names that ``help-or-hurt'' pattern as a
\emph{bottleneck signature} rather than as noise, and operationalizes routing
across heterogeneous interventions accordingly.

\subsection{Diagnostic evaluation of RAG}
Several frameworks now evaluate RAG without per-question gold answers.
RAGAS scores faithfulness, answer relevance, and context relevance using LLM
judges \citep{es2024ragas}; ARES trains lightweight evaluators against
synthetic judgments and includes prediction-powered confidence intervals
\citep{saadfalcon2024ares}; and RAGChecker decomposes errors along
retrieval-side and generation-side axes for fine-grained diagnosis
\citep{ru2024ragchecker}. A related thread documents that even well-retrieved
context is under-utilized when the relevant passage falls in the middle of a
long prompt \citep{liu2024lostmiddle}. These tools \emph{measure} RAG failure
modes but do not \emph{route} among interventions; our \controller{} closes
that loop by mapping the same kind of diagnostic signal onto a concrete
choice of policy per row.
